\section{Recovering the body and its quarter turns}\label{sec:recovery}

The caps in \cref{prop:equality-caps} determine the convex hull.
The constraints of the two original passages then put the original
body inside Romik's sofa. We finally use closedness to recover its
boundary as well as its area.

\begin{lemma}[Closed subsets of a regular closed set]\label{lem:regular-closed-recovery}
Let $F\subseteq\R^2$ be compact and regular closed, meaning
$F=\clos(\inter F)$. If a closed set $T$ satisfies $T\subseteq F$
and $|T|=|F|$, then $T=F$.
\end{lemma}

\begin{proof}
Finite equal area gives $|F\setminus T|=0$. The set
$\inter F\setminus T$ is open and has area zero, so it is empty:
every nonempty open subset of the plane contains a disk of positive
area. Thus $\inter F\subseteq T$. Closedness of $T$ gives
$F=\clos(\inter F)\subseteq T$.
\end{proof}

\begin{theorem}[Exact uniqueness]\label{thm:uniqueness}
For an ambidextrous sofa $S$,
\[
 |S|=\Mc\quad\Longleftrightarrow\quad
 S=R_\theta\SR+v\quad\text{for some }\theta\in\R,\ v\in\R^2.
\]
The equalities are equalities of compact sets.
\end{theorem}

\begin{proof}
Assume $|S|=\Mc$ and use \eqref{eq:original-equality-caps}.
Both actual caps have the same horizontal projection $I$.
Comparing the right endpoints of their translated reference
projections gives $r_{\KR}+a=r_{\KR}+b$, hence $a=b$.
Reflection commutes with horizontal translation. The hull identity
\eqref{eq:hull-from-caps} therefore gives
\begin{equation}\label{eq:original-equality-hull}
 K=U\cap\refl V
   =\bigl(\KR\cap\refl\KR\bigr)+(a,0)
   =\HR+(a,0).
\end{equation}
The final equality is the reference hull construction in
\cref{thm:reference}.

Translate horizontally by $-a$. Every point of $T-(a,0)$ is in
$\HR$. The full-angle passage constraints also impose
\[
 \nR(x)\le y\le1-\nR(x).
\]
Indeed the full-angle barriers of the translated caps are the
translated reference niche. By the definition of $\SR$ this proves
\begin{equation}\label{eq:reference-body-containment}
 T\subseteq\SR+(a,0).
\end{equation}
This step uses the constraints of the two original passages together
with their own aligned caps.

The reference body is regular closed. One can see the property
directly from its continuous niche and the strict bound
$\nR<1/2$. For $p=(x,y)\in\SR$, move vertically toward
$(x,1/2)$:
\[
 p_t=(x,(1-t)y+t/2),\qquad 0<t\le1.
\]
Reflection symmetry and convexity of $\HR$ put $(x,1/2)$ in $\HR$
whenever $p\in\HR$, so $p_t\in\HR$. Both niche inequalities at $p_t$ are
strict. The set on which these inequalities are strict is open.
Since the compact convex hull has nonempty interior,
$\HR=\clos(\inter\HR)$; every $p_t$ is consequently a limit of
interior hull points in that open set. Those points are in
$\inter\SR$. Sending $t$ to zero shows
\begin{equation}\label{eq:reference-regular-closed}
 \SR=\clos(\inter\SR).
\end{equation}


Now $T$ is compact, hence closed, and both sets in
\eqref{eq:reference-body-containment} have finite area $\Mc$.
Apply \cref{lem:regular-closed-recovery} to obtain
$T=\SR+(a,0)$. If the normalization was
$T=R_\eta S+v_0$, inversion gives the explicit equality
\[
 S=R_{-\eta}\SR+R_{-\eta}\bigl((a,0)-v_0\bigr).
\]


Conversely, every rotated and translated copy of $\SR$ is
ambidextrous: compose its two reference passages with the inverse
rigid map. Rigid maps preserve area, so every such copy has area
$\Mc$.
\end{proof}

\begin{theorem}[Quarter turns in every witnessing pair]\label{thm:quarter-turns}
Let $S$ be compact and connected, and let
$(\vartheta_-,v_-)$ and $(\vartheta_+,v_+)$ be any pair of
passages through $\Lm$ and $\Lp$ with the same initial rotation
map. If $|S|=\Mc$, then there are $t_-,t_+\in[0,1]$ such that
\[
 \vartheta_-(t_-)=\vartheta_-(0)-\frac\pi2,\qquad
 \vartheta_+(t_+)=\vartheta_+(0)+\frac\pi2.
\]
\end{theorem}

\begin{proof}
Apply the effective-angle construction to this particular pair of
passages, rebasing each lift at its own initial angle. Equality of
area forces both scalar scores to equal $\Mc/2$, so
\cref{thm:cap-equality} forces $\alpha=\gamma=\pi/2$.
For each motion the proper-terminal alternative in
\cref{lem:effective-angle} has angle strictly below $\pi/2$ and is
therefore excluded. Each original passage reaches its correctly
handed quarter turn.
\end{proof}

The theorem concerns angles reached during a passage. A witnessing
motion may continue rotating afterwards, reverse direction, or pause;
its terminal angle and its time parametrization are unrestricted by
this conclusion. Together, \cref{thm:optimality,thm:uniqueness,thm:quarter-turns}
prove \cref{thm:main}.
